\documentclass{article}

\PassOptionsToPackage{numbers, compress}{natbib}
% neurips_2026.sty is not in this tree and is not installed in the local TeX distribution
% (kpsewhich cannot find it), so requiring it makes the supplementary uncompilable on this
% machine. Use it when present -- the submission will need it -- and fall back to plain
% natbib otherwise so the document always builds while we are still filling in results.
\IfFileExists{neurips_2026.sty}{\usepackage[eandd]{neurips_2026}}{\usepackage{natbib}}

\usepackage[utf8]{inputenc}
\usepackage[T1]{fontenc}
\usepackage{hyperref}
\usepackage{url}
\usepackage{booktabs}
\usepackage{amsfonts}
\usepackage{amsmath}
\usepackage{nicefrac}
\usepackage{microtype}
\usepackage{graphicx}
\usepackage{xcolor}
\usepackage{placeins}
\newcommand{\todo}[1]{\textcolor{red}{\textbf{TODO:} #1}}

\title{Supplementary Material\\
\large When Alignment Hurts: A Controlled Benchmark for Ranking-Preserving 3D Face Metrics}

\author{%
  Anonymous Authors \\
  Anonymous Institution \\
  \texttt{anonymous@example.com}
}

\begin{document}
\maketitle

\noindent
This supplement reports experiments added after the first version of this work. Three
of them change how a result in the main paper should be read, and we state that
explicitly rather than only adding numbers: the per-mesh normalisation used to make
meshes comparable is not a neutral choice (Sec.~\ref{sec:supp_normalisation}); the
rigid-alignment baseline in the main comparison is not the strongest available one
(Sec.~\ref{sec:supp_m3dfb}); and the out-of-distribution transfer number reported for
FaceScape measured a change of topology-variant construction, not a change of morphable
model family (Sec.~\ref{sec:supp_transfer}).

%=====================================================================
\section{Per-Mesh Normalisation Is Not a Neutral Preprocessing Step}
\label{sec:supp_normalisation}

Every distance in this work is computed after a per-mesh normalisation that centres a
mesh on its vertex mean and divides by its largest absolute coordinate
(\texttt{maxabs}); this is the normalisation implemented by the dataset loader used for
both training and evaluation. Because the divisor is a single extreme coordinate, it is
itself topology-dependent: for one subject, the \texttt{original} and \texttt{down8k}
realisations differ by $4.5\%$ in scale and by $\approx0.04$ normalised units in
centroid position, which is the same order as the difference between two \emph{different}
identities.

We therefore recomputed raw Chamfer under a second, area-based normalisation --- centre
on the area-weighted centroid, divide by $\sqrt{A_{\mathrm{tot}}}$ --- on the identical
pair sets, and compared the two. Table~\ref{tab:supp_normalisation} groups the $30$
ordered topology pairs by the topologies they involve.

\begin{table}[!ht]
\centering
\small
\caption{Raw Chamfer under two per-mesh normalisations, Spearman correlation
($\uparrow$) with $D_{\mathrm{GT}}$ on the same pair sets ($30$ ordered topology pairs,
$13{,}050$ subject pairs, $30$ subjects). The learned metric is shown for reference. The
two normalisations fail on \emph{opposite} perturbations.}
\label{tab:supp_normalisation}
\begin{tabular}{lcccc}
\toprule
Topology pairs & $n$ & Chamfer (\texttt{maxabs}) & Chamfer (area) & Latent \\
\midrule
involving \texttt{crop}   & 10 & 0.087 & \textbf{0.279} & 0.612--0.725 \\
involving \texttt{noisy}  &  8 & \textbf{0.510} & 0.105 & 0.731--0.825 \\
clean topologies only     & 12 & 0.456 & \textbf{0.533} & --- \\
\midrule
all pairs                 & 30 & \textbf{0.253} & 0.172 & \textbf{0.730} \\
\bottomrule
\end{tabular}
\end{table}

Area normalisation improves $16$ of $30$ pairs and degrades $14$. The pattern is
mechanical rather than incidental. \texttt{maxabs} is set by one extreme vertex, and
\texttt{crop} removes exactly such vertices, so crop-involving pairs recover strongly
under area normalisation (\texttt{remesh}$\rightarrow$\texttt{crop}: $0.070
\rightarrow 0.451$). Area normalisation is set by total surface area, and \texttt{noisy}
inflates area, so noise-involving pairs collapse
(\texttt{noisy}$\rightarrow$\texttt{original}: $0.636 \rightarrow 0.066$).

Two consequences. First, part of the very low Chamfer correlations reported for
crop-involving pairs in the main paper is attributable to the normalisation rather than
to the distance itself, and we report both normalisations so that the geometric baseline
is credited with its better behaviour. Second, and more interesting for the argument of
this paper: normalisation is a \emph{second} mechanism, independent of pairwise
registration, by which a geometric distance loses the identity ordering when mesh
support changes. In aggregate \texttt{maxabs} ($0.253$) is the stronger of the two
choices, so the main paper's normalisation was not selected against the geometric
baseline; but no per-mesh normalisation is neutral, and the learned metric dominates
both in every group.

%=====================================================================
\section{Extended Baselines}
\label{sec:supp_baselines}

The first version compared the learned metric against raw Chamfer and registration
pipelines only. We add three further families, evaluated on the identical cross-topology
pair sets ($148{,}500$ pairs, $100$ subjects) so that the numbers are directly
comparable to the main tables.

\subsection{Correspondence-Free Geometric and Perceptual Distances}

\begin{table}[!ht]
\centering
\small
\caption{Additional baselines on clean cross-topology pairs, Spearman correlation
($\uparrow$) with $D_{\mathrm{GT}}$. Perceptual distances are cosine distances between
embeddings of canonical renders. LPIPS is reported on a $20$-subject subset because it
is pairwise in the image domain.}
\label{tab:supp_baselines}
\begin{tabular}{llc}
\toprule
Family & Metric & Spearman \\
\midrule
Learned 3D (ours)      & latent distance          & \textbf{0.751} \\
\midrule
Registration-free geom.& raw Chamfer              & 0.237 \\
                       & varifold (multi-scale)   & 0.151 \\
                       & currents (multi-scale)   & 0.074 \\
\midrule
Perceptual 2D          & ArcFace                  & 0.213 \\
                       & LPIPS ($20$ subj.)       & 0.120 \\
                       & CLIP ViT-B/32            & 0.105 \\
                       & DINOv2 ViT-S/14          & 0.074 \\
\bottomrule
\end{tabular}
\end{table}

\paragraph{Perceptual baselines are reported at their best, not their first, configuration.}
Meshes are rendered with a fixed orthographic camera and Lambertian shading. Our initial
renderer normalised each mesh independently, which framed the \texttt{crop} topology
larger than the other realisations of the same subject and therefore depressed exactly
the pairs that matter. We fixed this by deriving the camera centre and scale from each
subject's \texttt{original} mesh and reusing them for all of that subject's topologies,
and by averaging embeddings over three yaws ($-30^\circ, 0^\circ, +30^\circ$). This
raised ArcFace from $0.178$ to $0.213$, CLIP from $0.082$ to $0.105$ and DINOv2 from
$0.059$ to $0.074$. The gain is attributable to the shared frame and not to multi-view:
over crop-involving pairs ArcFace improves by $0.082$ on average against $0.023$
elsewhere, whereas in a single-topology control multi-view helps CLIP and \emph{harms}
DINOv2 ($0.402 \rightarrow 0.374$). Table~\ref{tab:supp_baselines} reports the corrected
numbers.

\paragraph{Currents do not rank identities, and this is informative.}
The currents distance uses oriented normals, which are as sensitive to tessellation as
to identity: over $8$ subjects, the ratio between the worst same-identity
cross-topology distance and the median cross-identity distance is $0.97$, i.e. the two
distributions overlap almost completely. We report currents as an orientation-sensitivity
diagnostic rather than as a ranking baseline. The varifold distance, which squares the
normal term and is therefore orientation-invariant, does rank identities (ratio $1.67$),
and is the number in Table~\ref{tab:supp_baselines}.

\subsection{M3DFB Error Estimators}
\label{sec:supp_m3dfb}

M3DFB~\cite{fg2025} provides a modular family of reconstruction-error estimators, built
as the product of a rigid aligner (ICP or landmark-based robust linear regression, RLR),
an optional non-rigid stage, and an optional topology corrector (ETC). We ran the
applicable estimators on our pair sets using the authors' implementation.

\begin{table}[!ht]
\centering
\small
\caption{M3DFB estimators on cross-topology pairs ($20$ subjects, $5{,}700$ pairs),
Spearman correlation ($\uparrow$) with $D_{\mathrm{GT}}$. \textbf{All M3DFB estimators
require 51 landmarks, which our setting does not have}; the landmarks supplied here are
exact on BFM-ordered topologies and nearest-vertex transferred on the others, so these
numbers are an upper bound on what the estimators would achieve on genuinely
landmark-free cross-topology data.}
\label{tab:supp_m3dfb}
\begin{tabular}{llc}
\toprule
& Estimator & Spearman \\
\midrule
& latent distance (ours) & \textbf{0.710} \\
\midrule
E9  & RLR + Chamfer corr. + P2P        & 0.456 \\
E10 & RLR + Chamfer corr. + P2P + ETC  & 0.417 \\
E2  & ICP + Chamfer corr. + P2P + ETC  & 0.290 \\
E1  & ICP + Chamfer corr. + P2P        & 0.284 \\
& raw Chamfer                          & 0.265 \\
\bottomrule
\end{tabular}
\end{table}

Three observations. (i) E1 reproduces the rigid-ICP-plus-Chamfer number of the main
paper ($0.284$ against $\approx0.29$), confirming that our own implementation of that
pipeline was faithful; the comparison now runs on the original authors' code.
(ii) \textbf{The landmark-based rigid aligner is far stronger than ICP} ($0.456$ against
$0.284$). ICP with scaling optimises away precisely the identity-related morphology the
ranking must preserve, whereas a landmark Procrustes fit preserves it. Reporting only
ICP would understate the registration-based family, and we therefore treat RLR as the
rigid baseline of record. (iii) The topology corrector barely affects cross-topology
ranking ($\mathrm{E2}-\mathrm{E1} = +0.006$, $\mathrm{E10}-\mathrm{E9} = -0.039$).
Even with landmarks supplied for free, the best M3DFB estimator remains $0.25$ Spearman
below the learned metric on the same pairs.

Twelve of the sixteen estimators are runnable but too slow for a benchmark of this size
(elastic linear regression alone costs $53$\,s per pair at $23$k vertices and scales
quadratically; non-rigid ICP costs $185$\,s per pair at $8$k vertices), and two of the
components are unusable on cross-topology data by construction --- identity
correspondence requires equal topology, and the landmark distance would score our own
landmark transfer.

%=====================================================================
\section{REMESH-2: A Second Morphable Model}
\label{sec:supp_remesh2}

To separate the effect of the morphable model family from the effect of mesh
representation, we extended the benchmark with identities from FLAME, generated by the
same protocol as the BFM half: $600$ identities sampled from the identity subspace with
truncated Gaussian coefficients, each realised in the same six topologies.

\paragraph{The face region is transported through correspondence, not redefined.}
A full FLAME head and a BFM face patch are not a fair cross-model pair. We therefore
carry the BFM crop region onto FLAME through the published BFM-to-FLAME vertex
correspondence: a FLAME vertex is retained if and only if its entire positional weight
falls inside the BFM crop index set. The criterion is sharp rather than tuned ($3022$
vertices at weight $\ge 0.99$, $3037$ at $\ge 0.25$), the two eyeball components absent
from the BFM patch are discarded, and the result is a fixed index set of $1930$ vertices
and $3770$ faces, identical for every identity.

\paragraph{Topology variants are defined by ratio, not by absolute count.}
The decimation and subdivision targets of the BFM pipeline ($16$k and $120$k triangles)
are calibrated on a $46{,}440$-triangle patch. Applied verbatim to FLAME's $3{,}770$
triangles both variants would be no-ops and the FLAME half of the benchmark would be
degenerate ($\texttt{down8k} \equiv \texttt{original}$). We therefore preserve each
model's ratio to its own \texttt{original}, which is the property that makes a variant
comparable across models; the realised ratios agree to three decimals
($0.344$ and $2.584$ on FLAME against $0.344$ and $2.587$ on BFM). The remaining
variants are defined by scale-relative rules and transfer unchanged.

\paragraph{Identity distinctness.}
Sampled identities could in principle be too close to be distinguishable. Over the $600$
FLAME identities the closest pair is at $30.8\%$ of the median pairwise distance, and the
minimum pairwise distance in coefficient space is $10.58$; no near-duplicate identities
are present.

%=====================================================================
\section{What the FaceScape Transfer Number Measured}
\label{sec:supp_transfer}

The first version of this work reported a cross-topology zero-shot Spearman correlation
of $0.115$ on FaceScape and concluded that transfer to a different mesh family is a
substantial domain shift. A controlled decomposition shows that this conclusion does not
follow from that experiment: the number is dominated by how the cross-topology
\emph{mate} of each pair was constructed, not by the change of morphable model.

All cells in Table~\ref{tab:supp_transfer} use the same released checkpoint, the same
evaluator and the same protocol; only the evaluation domain and the rule used to build
the second topology of each pair vary.

\begin{table}[!ht]
\centering
\small
\caption{Cross-model transfer of the REMESH-trained metric, Spearman correlation
($\uparrow$) with each domain's $D_{\mathrm{GT}}$. Rows differ only in the evaluation
domain and in the rule used to construct the cross-topology mate. The Poisson mate is
the construction used in the first version of this work.}
\label{tab:supp_transfer}
\begin{tabular}{llrc}
\toprule
Domain & Cross-topology mate & Pairs & Spearman \\
\midrule
BFM, held-out subjects        & benchmark variants        & 148{,}500 & \textbf{0.751} \\
\midrule
FLAME (ours)                  & remesh, benchmark rule    & 148{,}500 & 0.406 \\
FaceScape                     & remesh, benchmark rule    &  11{,}990 & 0.404 \\
FaceScape, support-matched    & remesh, benchmark rule    &  11{,}990 & 0.263 \\
\midrule
FLAME (ours)                  & Poisson reconstruction    & 148{,}500 & 0.109 \\
FaceScape                     & Poisson reconstruction    &  11{,}990 & 0.057 \\
\bottomrule
\end{tabular}
\end{table}

\paragraph{Under a matched pair construction, a different model family costs almost nothing.}
FaceScape reaches $0.404$ and our synthetic FLAME reaches $0.406$: statistically
indistinguishable, although FaceScape consists of real scans, of a different template,
of a different topology, and of a whole-face rather than cropped support. This is direct
evidence that the metric transfers across morphable model families without retraining or
adaptation.

\paragraph{The failure is surface reconstruction, and it is domain-independent.}
The mate used in the first version is a Poisson reconstruction at depth $7$, which closes
the boundary loops of the source mesh (minimum boundary radius $0.955$ against $0.549$),
adds $\approx21\%$ surface area, leaves the original support (maximum radius $1.412$
against $1.245$), and departs from its own source surface by a normalised Hausdorff
distance of $0.38$ --- against $0.088$ for the \texttt{original}/\texttt{remesh} pair the
metric was trained on. Applying that same construction to FLAME drops FLAME from $0.406$
to $0.109$, reproducing the apparent cross-model failure \emph{with no change of
morphable model at all}. The weakness is therefore robustness to surface
re-reconstruction, it is present in every domain including the training one, and it is a
separate axis from cross-model transfer. The Poisson mate remains a legitimate and
harder test; the error in the first version was attributing its outcome to the change of
model family.

\paragraph{Matching the support hurts, and that says something about the ground truth.}
Cropping FaceScape to a support matching the BFM patch \emph{lowers} the correlation from
$0.404$ to $0.263$. The likely reason is that $D_{\mathrm{GT}}$ is a per-vertex mean over
the whole face, so it integrates jaw, forehead and cheek geometry that the crop discards:
the mesh support and the support over which the ground truth is defined must agree. A
ground truth restricted to the cropped region would settle this, and requires dense
correspondence on the detail meshes, which the $10$k downsamples do not provide.

\paragraph{Which axes of domain shift matter.}
On normalised meshes, FaceScape is out of distribution relative to BFM on the region
axes (bounding-box $x/y$ $0.952$ against $0.773$; area within radius $0.6$ of the
centroid $0.247$ against $0.456$) and FLAME matches BFM there. On resolution, however,
FLAME is \emph{further} from BFM than FaceScape is (mean edge length over bounding-box
diagonal: $0.0175$ against $0.0054$ for BFM and $0.0085$ for FaceScape) while
transferring better, so resolution cannot be the causal axis. Non-dimensionalised, the
first Laplace--Beltrami eigenvalue agrees across all three domains
($\lambda_1 A \in \{11.1, 9.8, 8.2\}$), confirming that the pipeline is insensitive to
the raw unit scale of each domain: scaling a mesh's raw coordinates by $10^2$ changes the
resulting embedding by a relative $6\cdot10^{-4}$ (cosine similarity $1.000000$).

%=====================================================================
\section{Do Precomputed Operators Diffuse Heat Badly Across Topologies?}
\label{sec:supp_potential}

DiffusionNet learns diffusion times against a Laplace--Beltrami operator that is
\emph{precomputed per mesh}. Each realisation of an identity therefore carries its own
operator, and the network's invariance across topologies is only as good as the agreement
between those operators. This is not a concern for resampling --- the cotangent Laplacian
converges under refinement --- but it is a concern when the \emph{boundary} moves. A cropped
face is a manifold with a different boundary, and the Neumann condition implied by the
cotangent Laplacian reflects heat off that new boundary, so the same identity diffuses
differently depending on where it was cut. Our per-group numbers are consistent with this:
the crop group is the weakest of the three ($\rho = 0.763$ against $0.828$ for noise and
$0.816$ for pure resampling).

\paragraph{Intervention.} We test a boundary-insensitive operator built from the infinite
potential well of Liu, Jacobson and Crane~\cite{liu2017dirac}: replacing $L$ by
$L' = L + U$, where $U$ is a sigmoidal well on normalised geodesic distance from the mesh
centre,
\begin{equation}
U_i = m_i \, c \, \bigl(1 + e^{-\beta (d_i - \alpha)}\bigr)^{-1},
\qquad
c = c_{\mathrm{rel}} \max_k \lambda_k ,
\end{equation}
with $m_i$ the lumped mass, $d_i$ the heat-method geodesic distance from the area-weighted
centroid, $\beta = 100$ (the paper's $\gamma$) and $c_{\mathrm{rel}} = 10^{2}$. The well makes
the effective domain canonical: eigenfunctions are pushed away from the boundary, so a crop
that removes material outside the well leaves the retained spectrum nearly unchanged. Setting
$c$ \emph{relative} to $\lambda_k$ rather than to the paper's absolute $c = 10^{10}$ is a
deliberate deviation: the absolute value made the $8$k-decimation group $4.3\times$ worse
through mesh-dependent ill-conditioning.

\paragraph{Where the offset $\alpha$ is placed is the method.}
The source prescribes a \emph{collection-level} offset: ``$\beta$ is normalized such that
across an entire collection of patches no boundary points are contained in the region of
interest''~\cite{liu2017dirac}. This is what makes the construction work --- every patch is
restricted to the \emph{same} canonical domain, so their spectra are comparable. Our first
implementation instead took a per-mesh quantile of each mesh's own distance distribution,
which is not the same thing and defeats the mechanism: measured over meshes, it places the
transition at $\alpha \approx 0.68$--$0.73$ while the nearest boundary point sits at
$0.27$--$0.33$, so on $8$ of $8$ meshes inspected the boundary lay \emph{inside} the region
the well is meant to exclude, and the offset differed between topologies of the same identity
($0.679$ for \texttt{crop} against $0.727$ for \texttt{original}) --- exactly the shared
domain the well is supposed to establish. We therefore calibrate one offset and one common
length scale over the collection, taking $\alpha$ strictly below the smallest normalised
centre-to-boundary distance found anywhere ($\alpha = 0.210$, safety factor $0.9$), and report
the per-mesh variant as a placement ablation rather than as the method.

\paragraph{Pre-registered prediction, and its refutation.}
The mechanism is specific, so we stated the expected signature before measuring and report it
either way: \textbf{the crop group improves} and \textbf{the pure-resampling group is
unchanged}. Both halves are wrong. With the well placed at the offset our own spectral sweep
selects, crop \emph{falls} by $0.006$ and resampling \emph{falls} by $0.020$; the only group
that gains is \texttt{noisy} ($+0.009$), which we did not predict. Placing the well badly (a
per-mesh quantile) is not the explanation --- that arm behaves the same way, and the two
differ by less than the gap to the control.

\paragraph{A spectral proxy does not predict the metric.}
The more useful result is methodological. Table~\ref{tab:supp_alpha_sweep} measured, without
any training, that $\alpha = 0.55$ cuts cross-topology spectral inconsistency on the crop group
by $46\%$ while \emph{raising} identity separability --- both moving the right way at once, which
we read as the signature of a mechanism acting on the boundary. The trained metric then moved
the other way on that very group. Spectra that agree better across topologies therefore do not
imply identity orderings that survive better, and the cheap proxy would have chosen this
configuration on evidence that does not transfer.

We can say what the well costs. At $\alpha = 0.55$ the region of interest retains only
$22$--$50\%$ of the surface depending on the realisation ($22\%$ on \texttt{noisy}, whose extra
area falls outside the well; $46\%$ on \texttt{crop}). A network trained cross-topology has
already had to become robust to boundary changes, and it plausibly does so using exactly the
peripheral geometry the well removes; the invariance bought does not pay for the identity
information spent. This is a negative result about applying the construction in this regime, not
about the construction, whose behaviour on partial-shape spectra we reproduce
(Table~\ref{tab:supp_alpha_sweep}).

\paragraph{Protocol.} Both arms are trained here rather than reusing the v1 checkpoint, whose
subject subset and epoch budget differ. The arms share data, recipe, seed ($1234$), batch size
and epoch budget ($60$); they differ only in whether the cached operators come from $L$ or
$L'$. Model selection uses the same rule for both arms (best cross-topology mesh score), fixed
in advance --- selecting each arm by its own best group score would manufacture the effect
under test.

\paragraph{The prescribed offset is the worst choice on this data.}
Because the well trades boundary invariance against how much surface it retains, we measured
that trade-off directly on the spectrum, with no network involved: if the well does anything,
it must already show up in the eigenvalues. For each offset we report the mean spectral
distance between two \emph{topologies of the same identity} (which the well must reduce) and
between \emph{different identities in the same topology} (which a well that has eaten the face
also reduces). Reporting the first alone would make $\alpha \to 0$ look like a triumph, since
a small enough well turns every mesh into the same tiny disc. Table~\ref{tab:supp_alpha_sweep}
uses $40$ identities across the $6$ topologies.

\begin{table}[!ht]
\centering
\caption{Offset sweep, $40$ identities, no training. \texttt{within} is the mean spectral
distance between topologies of one identity, \texttt{between} that between identities within
one topology, and \texttt{crop-within} restricts \texttt{within} to pairs involving the
cropped realisation. The paper's collection-wide offset ($0.21$ here) is worse than using no
well at all on \emph{both} axes; the useful regime is an offset roughly $2.5\times$ larger.}
\label{tab:supp_alpha_sweep}
\begin{tabular}{lcccc}
\toprule
$\alpha$ & within & between & ratio $\uparrow$ & crop-within $\downarrow$ \\
\midrule
no well            & 0.2535 & 0.1774 & 0.700 & 0.3291 \\
$0.21$ (prescribed)& 0.5897 & 0.3752 & 0.636 & 0.4837 \\
$0.30$             & 0.4266 & 0.2912 & 0.683 & 0.3389 \\
$0.40$             & 0.2963 & 0.2195 & 0.741 & 0.2386 \\
$\mathbf{0.55}$    & \textbf{0.2208} & 0.1742 & \textbf{0.789} & \textbf{0.1765} \\
$0.75$             & 0.3218 & 0.2513 & 0.781 & 0.3376 \\
\bottomrule
\end{tabular}
\end{table}

At $\alpha = 0.55$ both quantities improve at once --- cross-topology inconsistency on the crop
group falls by $46\%$ ($0.3291 \rightarrow 0.1765$) while identity separability \emph{rises}
($0.700 \rightarrow 0.789$). That they move together, rather than trading off, is the signature
of a mechanism acting on the boundary rather than one simply discarding surface. The prescribed
$\alpha = 0.2098$ does the opposite on both counts.

The reason is a mismatch of regimes, not a defect of the construction. The condition ``no
boundary points inside the region of interest'' presumes patches cut from a larger surface,
whose shared interior is wide and whose boundary is far from the centre. A face mesh is not a
patch of anything: its boundary is intrinsic and close to the centre relative to its extent, so
enforcing that condition here retains only $2$--$7\%$ of the area. We therefore select $\alpha$
empirically and report the prescribed value as the endpoint of the sweep, rather than presenting
$0.55$ as if it followed from the source.

\begin{table}[!ht]
\centering
\caption{Ranking preservation (Spearman $\rho$ against $D_{\mathrm{GT}}$) by topology group,
$100$ held-out subjects, $148{,}500$ cross-subject pairs. All arms share data, recipe, seed and
epoch budget; each is read with the operators it was trained on; model selection uses one rule
throughout. Best per row in bold. \textbf{The potential well does not deliver what it was
predicted to deliver}: at the offset our own sweep selects it \emph{lowers} the crop group
($-0.006$) and costs $0.020$ on pure resampling, and only \texttt{noisy} improves ($+0.009$).}
\label{tab:supp_potential}
\begin{tabular}{lccc}
\toprule
 & no well & well misplaced & well $\alpha{=}0.55$ \\
Group & ($L$) & (per-mesh offset) & (placed) \\
\midrule
crop (boundary moved)     & \textbf{0.7072} & 0.7034 & 0.7012 \\
noisy                     & 0.7561 & 0.7573 & \textbf{0.7653} \\
resample (sampling only)  & \textbf{0.7719} & 0.7533 & 0.7521 \\
\midrule
all (pooled)              & \textbf{0.7347} & 0.7315 & 0.7215 \\
\bottomrule
\end{tabular}

\vspace{0.4em}
\footnotesize A fourth arm was designed and implemented, pooling restricted to the well's region
of interest so that the \emph{support} of the global descriptor would be corrected alongside the
operator. It is reported as absent rather than as pending: it never reached a compute node,
having sat in a shared queue carrying several hundred jobs for more than twenty hours, and the
three arms above already separate the well's presence from its placement, which is what the
table is for. Its code and its end-to-end verification are in the repository so the arm can be
run when capacity allows.
\end{table}

%=====================================================================
\section{Methodological Notes}
\label{sec:supp_methodology}

\paragraph{The ground truth depends on the space it is computed in.}
$D_{\mathrm{GT}}$ is defined by Eq.~(2) of the main paper, but the space in which the
vertices live is a choice. Computing it on the same FLAME meshes before and after the
per-mesh normalisation used by the evaluation gives two matrices whose Spearman
correlation is only $0.571$ (on BFM the corresponding figure is $0.81$). Two ground
truths built from the same meshes, differing only in normalisation, therefore order
identities substantially differently. We adopt the matrix computed in the normalised
space --- the space in which the evaluated metrics live --- as the reference, and report
the raw-space matrix as a robustness control.

\paragraph{The ground truth also depends on the topology it is measured in.}
The stronger version of the previous point is obtained by expressing the \emph{same}
identities in two topologies. Transporting our $500$ BFM identities through the
BFM-to-FLAME vertex correspondence yields, for each identity, a second mesh in FLAME
topology; $D_{\mathrm{GT}}$ is then computable in either topology by the same Eq.~(2),
since dense correspondence holds within each. Table~\ref{tab:supp_gt_topology} compares the
two resulting matrices over all $124{,}750$ identity pairs.

\begin{table}[!ht]
\centering
\small
\caption{Agreement between two correspondence-based ground truths built on the same $500$
identities, differing only in the topology in which the vertices are sampled. Chance levels
are $0.2\%$ for nearest-neighbour match and $1.0\%$ for top-5 overlap.}
\label{tab:supp_gt_topology}
\begin{tabular}{lccc}
\toprule
Space & Spearman & NN match & top-5 overlap \\
\midrule
raw                                & 0.865 & 34.8\% & 45.1\% \\
per-mesh normalised (as evaluated) & 0.573 & 16.4\% & 27.5\% \\
\bottomrule
\end{tabular}
\end{table}

The nearest-neighbour agreement is far above chance, so the identity structure does survive
re-topologisation; but agreement is well short of perfect, and in the space in which metrics
are actually evaluated the rank correlation between the two ground truths is only $0.573$.
The mechanism is the one identified in Sec.~\ref{sec:supp_transfer}: FLAME topology samples
the same surface with $1930$ vertices rather than $23{,}470$, so a per-vertex mean weights
facial regions differently.

Consequently $D_{\mathrm{GT}}$ should be read as \emph{defined relative to a reference
topology and a normalisation space}, not as an absolute quantity. A metric evaluated on
pairs that cross topologies is scored against a target that itself moves when the reference
topology changes, and a benchmark of this kind must declare both. This is the thesis of this
paper taken one step further: it is not only registration-based pipelines that lose the
identity ordering under a change of mesh representation --- a correspondence-based ground
truth partially loses it too, once no topology is privileged.

\paragraph{Cross-domain identity distances are undefined, not zero.}
There is no identity distance between a BFM subject and a FLAME subject: the two
identity spaces are unrelated. A joint ground truth over both domains is therefore
block-diagonal with \emph{undefined} off-diagonal blocks. Any joint training must never
read those entries; we enforce this by forming every mini-batch from a single domain, so
that all pairs within a batch are intra-domain, and by storing the off-block entries as
\texttt{NaN} so that an error surfaces rather than being silently trained on.

\paragraph{Batch size is a property of the objective, not only of throughput.}
The stress and ranking terms are computed over pairs within a batch, so the number of
supervised pairs grows quadratically in the number of subjects per batch: the
configuration of the main paper ($5$ subjects) provides $10$ subject pairs, from which
the hard-pair mining fraction of $0.7$ has little to select. At the extreme, a
$2$-subject batch makes the subject-level stress term identically zero, because both
distance matrices are normalised by their own off-diagonal mean. We report batch size as
an ablation rather than fixing it silently.

\paragraph{Operators are recomputed per realisation.}
For completeness: the DiffusionNet operators are computed on each topology realisation
separately, not inherited from the \texttt{original} mesh. The \texttt{noisy} variant is
the tightest check, since it has the same vertex count as \texttt{original} but different
geometry; its Laplacian values differ, as they must.

{\small
\bibliographystyle{IEEEtran}
\bibliography{references}
}

\end{document}
